\section{$k$-Median Algorithm Opening $k + O(\log n)$ Centers}
\label{sec:merging}

In this section we expand on the ideas in Section~\ref{sec:logad}, based on the techniques developed in~\cite{MR4928455-Cohen25}, in order to obtain an algorithm for $k$-median that opens at most $k + O(\log n)$ centers.
As in their paper, this requires us to develop a robust version of the Log-Adaptive Algorithm, and to slightly adapt their merging procedure to our setting.
We remark that the definitions and notation in this section follow closely those in~\cite{MR4928455-Cohen25}.

First, we generalize the notions of $1$-openability and $2$-openability in Definitions~\ref{def:openable-mod-1} and~\ref{def:openable-mod-2} to approximate notions that we call $(1,\eta)$-openability and $(2,\eta)$-openability.
Note that these definitions of openable only depend on $\theta$ and $S$.
We will think of $\eta$ throughout the analysis as $2^{-n \cdot |F|}$.

\begin{definition}[$(1, \eta)$-openable facility]
\label{def:openable-mod-eta-1}
Consider a time $\theta$ and a set $S$ of open facilities.
Let $I = \{j \in D \mid \theta \geq (1-\delta) d(j, S)\}$ and $A = D \setminus I$.
A facility $i\in F\setminus S$ is $(1,\eta)$-\emph{openable} w.r.t.\ $(S,\theta)$ if
satisfying:
\[
\sum_{j\in A}\beta(\theta, (1-\delta)\dist(i,j))
\;+\;
\sum_{j\in I}\beta((1-\delta)\dist(j,S), (1-\delta)\dist(i,j))
\;\ge\; \hat f - n \eta.
\]
\end{definition}

\begin{definition}[$(2, \eta)$-openable facility]
\label{def:openable-mod-eta-2}
Consider a time $\theta$ and a set $S$ of open facilities.
Let $I = \{j \in D \mid \theta \geq (1-\delta) d(j, S)\}$ and $A = D \setminus I$.
A facility $i\in F\setminus S$ is $(2,\eta)$-\emph{openable} w.r.t.\ $(S,\theta)$ if there exist bids $(\tau_j)_{j\in A}$
satisfying:
\begin{enumerate}
\item[\textbf{(i)}] (\textbf{Only nearby clients may increase}) 
\[
\tau_j=\theta \qquad \forall j\in A\setminus \Ball{i}{\epsilon\theta}.
\]
\item[\textbf{(ii)}] (\textbf{Nearby clients increase only slightly})
\[
\theta \le \tau_j \leq \min\Bigl\{(1-\delta)\dist(j,S),\, (1+\epsilon^2) \theta \Bigr\}
\qquad \forall j\in A\cap \Ball{i}{\epsilon\theta}.
\]
\item[\textbf{(iii)}] (\textbf{Facility $i$ is paid for in the $(1-\delta/2)$-scaled metric})
\[
\sum_{j\in A}\beta(\tau_j, (1-\delta/2)\dist(i,j))
\;+\;
\sum_{j\in I}\beta((1-\delta)\dist(j,S), (1-\delta/2)\dist(i,j))
\;\ge\; \hat f - n \eta.
\]
\item[\textbf{(iv)}] (\textbf{No overbidding for regular facilities after updating to $\tau$})
\[
\forall i_0\in F:\qquad
\sum_{j\in A}\beta(\tau_j,\dist(i_0,j))
\;+\;
\sum_{j\in I}\beta((1-\delta)  \dist(j,S), \dist(i_0,j))
\;\le\; \hat f.
\]
\item[\textbf{(v)}] (\textbf{Triangle inequality})
\[
\tau_j \le \tau_k + \dist(j,k)
\qquad \forall j,k\in A.
\]
\end{enumerate}
\end{definition}

As in \cite{MR4928455-Cohen25}, we also allow the introduction of ``free facilities'' whose opening costs are not paid.
Let $S_{\mathrm{free}} \subseteq S$ be the multiset of free facilities that have been opened.
For each free copy $\tilde{i} \in S_{\mathrm{free}}$, we also define a parameter $u(\tilde{i})$ which represents the distance between the free facility $\tilde{i}$ and its original copy $i$.
For any point $x$ in the metric space, the distance is given by $d(\tilde{i}, x) = d(i, x) + u(\tilde{i})$.
Also let $S_{\mathrm{reg}} = S \setminus S_{\mathrm{free}}$ be the set of open regular facilities.

For $\tau \in \{1, 2\}$, we define $(\nu,\eta)$-valid sequences.
\begin{definition}[$(\nu, \eta)$-valid sequence]
\label{def:valid-sequence-eta}
Consider a time $\theta$ and let $S$ be a set of opened facilities. 
A sequence $(h_1, \ldots, h_{\ell})$ of facilities is $(\nu,\eta)$-valid if the following conditions hold.
\begin{itemize}
    \item \textbf{(Validity)} Each $h_t$ is either free or $(\nu,\eta)$-openable w.r.t. the time $\theta$ and a super set $S' \supseteq S \cup \{h_1, \ldots, h_{t-1}\}$.
    \item \textbf{(Maximality)} The sequence is maximal, i.e., there is no $(\nu,0)$-openable facility w.r.t. the time $\theta$ and $S \cup \{h_1, \ldots, h_{\ell}\}$.
\end{itemize}
\end{definition}

We note that Algorithm~\ref{alg:log-adaptive-greedy} generates $(\nu,0)$-valid sequences:

\begin{lemma}
\label{lem:log-valid-1}
Any sequence generated in Stage 1 of Algorithm~\ref{alg:log-adaptive-greedy} is $(1,0)$-valid.
\end{lemma}
\begin{proof}
Both properties follow from the definition of the algorithm.
For the first property, a facility $h_t$ is added only if it is $(1,0)$-openable w.r.t. the time $\theta$ and $S \cup \{h_1, \ldots, h_{t - 1}\}$.
For the second property, the loop in Stage 1 continues as long as there exists a $(1,0)$-openable facility w.r.t. $\theta$ and the set of facilities opened so far.
\end{proof}

\begin{lemma}
\label{lem:log-valid-2}
Any sequence generated in Stage 2 of Algorithm~\ref{alg:log-adaptive-greedy} is $(2,0)$-valid.
\end{lemma}
\begin{proof}
Both properties follow from the definition of the algorithm.
For the first property, a facility $h_t$ is added only if it is $(2,0)$-openable w.r.t. the time $\theta$ and $S \cup \{h_1, \ldots, h_{t - 1}\}$.
For the second property, the loop in Stage 2 continues as long as there exists a $(2,0)$-openable facility w.r.t. $\theta$ and the set of facilities opened so far.
\end{proof}

For simplicity of notation, let $\nu : \mathbb{N} \to \{1,2\}$ satisfy $\nu(x) = 1$ for odd integers and $\nu(x)=2$ for even integers.
We define now an $\eta$-valid sequence of sequences, which consists of $(\nu,\eta)$-valid sequences and in addition the first sequence is required to be of type $1$ and the sequences are required to alternate between type $1$ and type $2$.

\begin{definition}
\label{def:valid-seq-of-seq}
For a sequence of sequences $\mathcal{H} = (H_1, \ldots, H_L)$, we say it consists of $\eta$-valid sequences if $H_{t}$ is $(\nu(t), \eta)$-valid w.r.t. $\theta = (1+\epsilon^2)^{\lfloor (t-1)/2 \rfloor}$ and $S = \cup_{q \in [t - 1]} H_q$, for all $1 \leq t \leq L$.
\end{definition}
Furthermore, we say that $\mathcal{H}$ is a solution if $L$ is even and $A = \emptyset$ at the end, meaning that at time $\theta=(1+\epsilon^2)^{|\mathcal{H}| / 2}$, the set of inactive clients $I = \{j \in D \mid \theta \geq (1-\delta) d(j, S)\}$ is equal to the set of all clients.
Otherwise, we say that $\mathcal{H}$ is a partial solution.

We prove that the following approximation guarantee holds for solutions $\mathcal{H}$:

\begin{theorem}
\label{theorem:alg-lmp-robust}
Let $\gamma \geq 1$ and let $\beta(\alpha,d) = [\alpha - \gamma \cdot d]^+$. Let $\rho \geq 1.01 \gamma$ be an upper bound on the optimal value of the factor-revealing program in Definition~\ref{def:frp} over all $n \geq 1$.
Let $\delta > 0$ be small enough compared to $1/\gamma$ and let $\eps > 0$ be small enough compared to $1/\gamma$ and $\delta$,
and let $\mathrm{OPT}$ be the optimal cost to the facility location problem with facility cost $f = \frac{\hat{f}}{\rho + 6 \delta}$ for $\hat f > n \eta$.
If $\mathcal{H}$ is an $\eta$-valid solution with respect to these parameters, then the set $S^\star$ of facilities opened by $\mathcal{H}$ satisfies
\[\sum_{j \in D} d(j, S^\star) \leq (1+8\delta) \rho \cdot (\mathrm{OPT} - |S^\star_{\mathrm{reg}}| \cdot f) + 2 |S^\star_{\mathrm{reg}}| \cdot n \eta\,.\]
\end{theorem}

We will prove Theorem~\ref{theorem:alg-lmp-robust} in Section~\ref{sec:logad-robust}.
The proof of the theorem follows the same set of steps as that of Theorem~\ref{theorem:alg-lmp}, but becomes more technically intricate because of the more general definition of $(\nu,\eta)$-openable.
An important technical challenge comes from the fact that a facility may be opened with respect to a \emph{superset} of the set of open facilities; this ends up making it more challenging to embed a solution into a feasible solution for the factor-revealing program, and requires us to use a more complicated assignment of values to the variables $r$ in the factor-revealing program.

\subsection{Merging Procedure}

Next, we will start with some definitions required in the merging procedure.

We define now two procedures $\mathrm{CompleteSolution}_{\hat f,u}$ and $\mathrm{CompleteSequence}_{\hat f,u}$, that are extensively used in the merging algorithm.
These procedures are defined analogously to the same procedures in~\cite{MR4928455-Cohen25}:

\begin{itemize}
    \item $\mathrm{CompleteSolution}_{\hat f,u}(\mathcal{H})$ takes a partial $\eta$-valid solution $\mathcal{H}$ as input and returns a solution by running Algorithm~\ref{alg:log-adaptive-greedy} using opening cost $\hat f$ and distance $u$ for free facilities already opened.
    If $|\mathcal{H}|$ is odd, it starts at Stage 2 with $\theta = (1+\eps^2)^{\lfloor (|\mathcal{H}| - 1) / 2 \rfloor}$, and if $|\mathcal{H}|$ is even, it starts at Stage 1 with $\theta = (1+\eps^2)^{|\mathcal{H}| / 2}$.
    Note that $\mathrm{CompleteSolution}$ can only terminate after a Stage 3 in Algorithm~\ref{alg:log-adaptive-greedy}, so the returned solution has even length.
    \item $\mathrm{CompleteSequence}_{\hat f,u}(\mathcal{H}=(H_1, \ldots, H_{p-1}),H_p = (h_1, \ldots, h_\ell))$ where $\mathcal{H}$ is a partial $\eta$-valid solution and $H_p$ is a $(\nu(p), \eta)$-valid sequence, except that it may not be maximal.
    If $p$ is odd, $\mathrm{CompleteSequence}$ runs Stage 1 in Algorithm~\ref{alg:log-adaptive-greedy} using opening cost $\hat f$ and distance $u$ for free facilities, and if $p$ is even it runs Stage 2 in Algorithm~\ref{alg:log-adaptive-greedy} using opening cost $\hat f$ and distance $u$ for free facilities, to return a maximal sequence that contains $H_p$ as a prefix.
\end{itemize}

The definitions of $\mathrm{CompleteSolution}$ and $\mathrm{CompleteSequence}$ imply that the solutions and the sequences generated by them are valid:

\begin{observation}
\label{obs:alg-valid}
Any output of $\mathrm{CompleteSolution}_{\hat f,u}(\mathcal{H})$ on a valid input is an $\eta$-valid solution,
and any output of $\mathrm{CompleteSequence}_{\hat f,u}(\mathcal{H}=(H_1, \ldots, H_{p-1}),H_p = (h_1, \ldots, h_\ell))$ on a valid input produces a $(\nu(p),\eta)$-valid partial sequence that contain $H_p$ as a prefix.
\end{observation}

We will also need a non-overbidding property to be satisfied at the beginning of each sequence in a solution generated by $\mathrm{CompleteSolution}$:

\begin{lemma}
\label{lem:seq-overbidding}
Let $\mathcal{H} = (H_1, \ldots, H_{p})$ be an $\eta$-valid partial solution.
Suppose that for $S' = \bigcup_{H \in \mathcal{H}} H$, $I' = \{j \in D \mid \theta \geq (1-\delta) d(j, S')\}$, and $A' = D \setminus I'$ it holds for all $i_0 \in F$ that
    \[\sum_{j\in A'}\beta(\theta,\dist(i_0,j))
\;+\;
\sum_{j\in I'}\beta((1-\delta)  \dist(j,S'), \dist(i_0,j))
\;\le\; \hat f.\]
Then, for all $H_t$ with $t > p$ in the output of $\mathrm{CompleteSolution}(\mathcal{H})_{\hat f, u}$, for $S = S' \cup \{H_{p+1}, \ldots, H_{t}\}$, $I = \{j \in D \mid \theta \geq (1-\delta) d(j, S)\}$, and $A = D \setminus I$, we have for all $i_0 \in F$ that 
    \[\sum_{j\in A}\beta(\theta,\dist(i_0,j))
\;+\;
\sum_{j\in I}\beta((1-\delta)  \dist(j,S), \dist(i_0,j))
\;\le\; \hat f.\]
\end{lemma}
\begin{proof}
The claim is captured by Lemma~\ref{lem:overbidding-robust}, which proves that the property holds even after the opening of each facility in each sequence.
\end{proof}

We now discuss the merging procedure itself.
Informally, the merging procedure takes two $\eta$-valid solutions $\mathcal{H}$ and $\mathcal{H}'$ that open more and less than $k$ regular facilities, respectively, and merges them to obtain an $\eta$-valid solution that opens exactly $k$ regular facilities.
This merging procedure is essentially identical to that in~\cite{MR4928455-Cohen25}, except that (1) it uses \emph{our} definitions of $\completesol$ and $\completesequence$ above and (2) it has a slightly different initialization.

Let us describe our initialization.
We perform bisection search on $\hat f$ in the interval $[1/n^2, 4\gamma n\Maxdist]$ and run $\completesol_{\hat f, u}(\emptyset)$ with empty $u$.
Let us start by analyzing the endpoints of this interval:
\begin{itemize}
    \item $\hat f = 1 /n^2$: At $\hat f = 1/n^2$ every client $j \in D$ satisfies $d(j, S^\star) \leq \frac{1+\eps^2}{1-\delta} (\gamma d(j, F) + 1/n^2)$ because each client may open its closest facility by itself by the time $\theta \geq \gamma d(j, F) + 1/n^2$.
    Then the resulting solution is a $\frac{\gamma+O(\eps^2)+O(1/n)}{1-\delta}$-approximation, which by our choice of parameters is also a $\rho$-approximation.
    If this solution opens at most $k$ centers, we are done.
    Else, it gives a solution that opens more than $k$ centers.
    \item $\hat f = 4 \gamma n \Maxdist$: At $\hat f = 4 \gamma n \Maxdist$ the algorithm opens exactly one facility, because when first facility is opened then every other client becomes immediately inactive at that step.
    If $k=1$, we are done.
    Else, it gives a solution that opens less than $k$ centers.
\end{itemize}
Then we perform bisection search until we find two $\hat f > \hat f'$ such that $|\hat f - \hat f'| \leq \eta$ and the corresponding solutions $\mathcal{H}$ and $\mathcal{H}'$ open strictly more and strictly less than $k$ facilities, respectively.

Then, as in~\cite{MR4928455-Cohen25}, the procedure maintains two solutions $\mathcal{H}$ and $\mathcal{H}'$ whose number of opened regular facilities sandwich $k$.
It also ensures that $\mathcal{H}$ is $\eta$-valid w.r.t. some $\hat f, u$ and $\mathcal{H}'$ is $\eta$-valid w.r.t. some $\hat f', u'$ such that $\hat f, u$ and $\hat f', u'$ differ in at most one ``entry'', by at most $\eta$: either they are identical, or $\hat f \neq \hat f'$ and $|\hat f - \hat f'| \leq \eta$, or there exists a single free facility $\tilde{i}$ with $u(\tilde i) \neq u'(\tilde i)$  and $|u(\tilde i) - u'(\tilde i)| \leq \eta$.
The procedure modifies $\mathcal{H}$ and $\mathcal{H}'$ in steps, such that after that $p$-th step of the procedure either it has found a solution that opens exactly $k$ facilities or it ensures that the first $p$ sequences of $\mathcal{H}$ and $\mathcal{H}'$ are identical.
Also, each step of the procedure opens at most three new free facilities.

For $\hat f \leq 4\gamma n \Maxdist$, at the latest when $\theta$ reaches $6 \gamma \Maxdist$ a facility will open and every client will become inactive, so $\theta$ is upper bounded by $(1+\eps^2) 6 \gamma \Maxdist$.
Because $\theta$ starts at $1$ and increases exponentially in the number of sequences (specifically, it is multiplied by $(1+\eps^2)$ every two sequences), this implies that the number of sequences in any $\eta$-valid solution with $\hat f \leq 4 \gamma n \Maxdist \leq \poly(n)$ is $O(\log n / \eps^2)$.
This implies that the procedure finds a solution within $O(\log n / \eps^2)$ steps, because otherwise it would reach the contradiction that $\mathcal{H}$ and $\mathcal{H}'$ are identical but open different numbers of regular facilities.

\begin{proof}[Proof sketch of Theorem~\ref{thm:k-median}]
The merging procedure generates an $\eta$-valid solution that opens exactly $k$ regular facilities and at most $O(\log n/ \eps^2)$ free facilities. 
Then, by Theorem~\ref{theorem:alg-lmp-robust},
\begin{align*}
\sum_{j \in D} d(j, S^\star)
&\leq (1+8\delta) \rho \cdot (\mathrm{OPT}_f - |S^\star_{\mathrm{reg}}| \cdot f) + 2 |\Sr^\star| \cdot n \eta\\
&= (1+8\delta) \rho \cdot \left( \mathrm{OPT}_f - f \cdot k \right)  + 2 |\Sr^\star| \cdot n \eta\\
&\leq \Paren{(1+8\delta)\rho + 2 |\Sr^\star| \cdot n \eta} \cdot \sopt\\
&\leq (1+9\delta)\rho \cdot \sopt\,,
\end{align*}
where $\sopt$ is the cost of the optimal $k$-median solution.
We used above that $\sopt \geq \mathrm{OPT}_f - f \cdot k$, that $\sopt \geq 1$ (because every distance is at least $1$), and also that $\eta = 2^{-n \cdot |F|}$.
The statement of Theorem~\ref{thm:k-median} is captured by renaming $\delta \gets 9\delta$ and noticing that we only choose $\eps$ to be a constant factor smaller than $\delta$.

Strictly speaking, $S^\star$ may not be a valid solution to the $k$-median problem because it also contains free facilities that are not members of $F$.
However, each free facility $\tilde{i}$ is a copy of a regular facility $i \in F$, so we consider in fact the solution that for each free facility $\tilde{i} \in S^\star$ opens the regular facility $i$ that it is a copy of.
This cannot increase the cost of the solution: for all $j \in D$, we have $d(j, i) = d(j, \tilde{i}) - u(\tilde{i}) \leq d(j, \tilde{i})$.
Hence, this solution also has cost upper bounded by $(1+9\delta)\rho \cdot \sopt$.
The number of regular facilities opened is still bounded by $k + O(\log n)$.
\end{proof}

\subsection{Merging Procedure Analysis}

The analysis of the merging procedure depends on the following two lemmas that are essential to argue that, during the merging procedure, the sequences produced are maximal.
An analogous result is implicit in~\cite{MR4928455-Cohen25}.

\begin{lemma}
\label{lem:pay-to-open-1}
Suppose $\hat f > 0$.
Let $H_p = (h_1, \ldots, h_\ell)$ be a $(1,\eta)$-valid sequence with respect to time $\theta$ and open facilities $S'$, except that it may not be maximal.
For $I' = \{j \in D \mid \theta \geq (1-\delta) d(j, S')\}$, $A' = D \setminus I'$, suppose that for all $i_0 \in F$,
    \[\sum_{j\in A'}\beta(\theta,\dist(i_0,j))
\;+\;
\sum_{j\in I'}\beta((1-\delta)  \dist(j,S'), \dist(i_0,j))
\;\le\; \hat f.\]
Let $S = S' \cup \{h_1, \ldots, h_{\ell}\}$, and let $I = \{j \in D \mid \theta \geq (1-\delta) d(j, S)\}$ and $A = D \setminus I$.
Suppose furthermore that
\[\sum_{j \in A} \beta(\theta, (1-\delta) d(i, j)) + \sum_{j \in I}\beta((1-\delta)d(j,S), (1-\delta)d(j, i)) \geq \hat{f}\]
for some $i \in F$. Then there exists some $i^\star \in F \setminus S$ that is $(1,0)$-openable with respect to $\theta$ and $S$.
\end{lemma}
\begin{proof}
We start by noting that $i \in S$ is impossible unless $f=0$, because then $\theta \leq (1-\delta) d(j,S) \leq (1-\delta)d(i,j)$ for all $j \in A$ and $d(j, S) \leq d(i,j)$ for all $j \in [n]$,  which implies that the left-hand side of the budget assumption is zero.

Then we argue that $i$ itself is $(1, 0)$-openable.
The first constraint follows from the given assumption.
Also, the third constraint is trivial because $\alpha_j = \theta$ for all $j \in A$.

For the second constraint, we start by observing that
\[\sum_{j \in A'} \beta(\theta, d(i_0, j)) + \sum_{j \in I'}\beta((1-\delta)d(j,S'), d(i_0, j)) \leq \hat{f}\,,\]
so because $d(j,S) \leq d(j, S')$ also 
\[\sum_{j \in A'} \beta(\theta, d(i_0, j)) + \sum_{j \in I'}\beta((1-\delta)d(j,S), d(i_0, j)) \leq \hat{f}\,.\]
Let us now consider the clients $j \in A' \cap I$.
Because $j \in I$, we have $\theta \geq (1-\delta) d(j,S)$, so we also have $\beta(\theta, d(i_0, j)) \geq \beta((1-\delta)d(j,S), d(i_0, j))$. 
Hence, overall, we also have
\[\sum_{j \in A} \beta(\theta, d(i_0, j)) + \sum_{j \in I}\beta((1-\delta)d(j,S), d(i_0, j)) \leq \hat{f}\,.\]
\end{proof}

\begin{lemma}
\label{lem:pay-to-open-2}
Suppose $\hat f > 0$.
Let $H_p = (h_1, \ldots, h_\ell)$ be a $(2,\eta)$-valid sequence with respect to time $\theta$ and open facilities $S'$, except that it may not be maximal.
For $I' = \{j \in D \mid \theta \geq (1-\delta) d(j, S')\}$, $A' = D \setminus I'$, suppose that for all $i_0 \in F$,
    \[\sum_{j\in A'}\beta(\theta,\dist(i_0,j))
\;+\;
\sum_{j\in I'}\beta((1-\delta)  \dist(j,S'), \dist(i_0,j))
\;\le\; \hat f.\]
Let $S = S' \cup \{h_1, \ldots, h_{\ell}\}$, and let $I = \{j \in D \mid \theta \geq (1-\delta) d(j, S)\}$ and $A = D \setminus I$.
Suppose furthermore that there exist values $(\tau_j)_{j \in A}$ such that $\tau_j = \theta$ for $j \in A \setminus B(i, \epsilon \theta)$, $\theta \leq \tau_j \leq \min((1-\delta)d(j,S), (1+\epsilon^2)\theta)$ for $j \in A \cap B(i, \epsilon \theta)$, $\tau_j \leq \tau_k + d(j, k)$ for all $j, k \in A$, and
\[\sum_{j \in A} \beta(\tau_j, (1-\delta/2) d(i, j)) + \sum_{j \in I}\beta((1-\delta)d(j,S), (1-\delta/2)d(j, i)) \geq \hat{f}\]
for some $i \in F$. Then there exists some $i^\star \in F \setminus S$ that is $(2,0)$-openable with respect to $\theta$ and $S$.
\end{lemma}
\begin{proof}
We start by noting that $i \in S$ is impossible unless $\hat f=0$, because then $\tau_j \leq  (1-\delta) d(i,j)$ for all $j \in A$ and $d(j, S) \leq d(j, i)$ for all $j \in [n]$, which implies that the left-hand side of the budget assumption is zero.

Fix $(\tau_j)_{j \in A}$ such that $\tau_{\max} = \max_{j \in A} \tau_j$ is minimized and the assumptions are satisfied with respect to some $i \in F \setminus S$.
We will prove that $i$ is $(2,0)$-openable.
The first three constraints and the fifth constraint are satisfied by assumption.

\textbf{Fourth constraint for $i_0 \in F \setminus S$.}
For the fourth constraint, suppose for the sake of contradiction that there exists some $i_0 \in F \setminus S$ such that
\[\sum_{j \in A} \beta(\tau_j, d(i_0, j)) + \sum_{j \in I}\beta((1-\delta)d(j,S), d(i_0, j)) > \hat{f}\,.\]

\textbf{Existence of non-zero active bid.}
Let us prove that there exists $j \in A$ such that $\tau_j > \theta$ and $\beta(\tau_{j}, d(i_0, {j})) > 0$.
We have that
\[\sum_{j \in A'} \beta(\theta, d(i_0, j)) + \sum_{j \in I'}\beta((1-\delta)d(j,S'), d(i_0, j)) \leq \hat{f}\,,\]
so because $d(j, S) \leq d(j, S')$ also 
\[\sum_{j \in A'} \beta(\theta, d(i_0, j)) + \sum_{j \in I'}\beta((1-\delta)d(j,S), d(i_0, j)) \leq \hat{f}\,.\]
Let us now consider the clients $j \in A' \cap I$.
Because $j \in I$, we have $\theta \geq (1-\delta) d(j,S)$, so we also have $\beta(\theta, d(i_0, j)) \geq \beta((1-\delta)d(j,S), d(i_0, j))$.
Hence, overall, we also have that 
\[\sum_{j \in A} \beta(\theta, d(i_0, j)) + \sum_{j \in I}\beta((1-\delta)d(j,S), d(i_0, j)) \leq \hat{f}\,.\]
Then, for this inequality to be violated when switching from $\theta$ to $(\tau_j)_{j \in A}$, there must exist some $j \in A$ such that $\tau_{j} > \alpha_{j}$ and $\beta(\tau_{j}, d(i_0, {j})) > 0$.

\textbf{Finalizing fourth constraint for $i_0 \in F \setminus S$.}
For some $\zeta > 0$ small enough, consider a new set of values $\{\tau'_j\}_{j \in A}$ such that $\tau'_j = \tau_j - \zeta$ if $\tau_j > \theta$ and $\tau'_j = \tau_j$ otherwise.
Then we still have $\tau'_j \geq \theta$ for all $j \in A$ and 
\[\sum_{j \in A} \beta(\tau'_j, d(i_0, j)) + \sum_{j \in I}\beta((1-\delta)d(j,S), d(i_0, j)) > \hat{f}\,.\]
Let us prove now that this leads to a contradiction.
Let $\tau'_{\max} = \max_{j \in A} \tau'_j$. We have $\theta < \tau'_{\max} < \tau_{\max}$.
Let $\alpha'_j = \min(\tau'_{\max}, (1-\delta)d(j,S))$, and note that $\alpha'_j \geq \tau'_{j}$ for all $j \in A$, because if $\alpha'_j = (1-\delta) d(j,S)$ we have $(1-\delta) d(j,S) \geq \tau_j \geq \tau'_j$.
Then we also have
\[\sum_{j \in A} \beta(\alpha'_j, d(i_0, j)) + \sum_{j \in I}\beta((1-\delta)d(j,S), d(i_0, j)) > \hat{f}\,.\]
Let $\tilde{\tau}_j = \min(\alpha'_j, (1+\epsilon) \theta - d(i_0, j))$ for $j \in A \cap B(i_0, \epsilon \theta)$ and $\tilde{\tau}_j = \theta$ for $j \in A \setminus B(i_0, \epsilon \theta)$.
We claim that 
\begin{align*}
&\sum_{j \in A} \beta(\tilde\tau_j, (1-\delta/2) d(i_0, j)) + \sum_{j \in I}\beta((1-\delta)d(j,S), (1-\delta/2)d(j, i_0))\\
=& \sum_{j \in A \cap B(i_0, \eps\theta)}  \beta(\min(\alpha'_j, (1+\epsilon) \theta - d(i_0, j)), (1-\delta/2) d(i_0, j)) + \sum_{j \in A \setminus B(i_0, \eps\theta)}  \beta(\theta, (1-\delta/2) d(i_0, j))\\
&\quad + \sum_{j \in I}\beta((1-\delta)d(j,S), (1-\delta/2)d(j, i_0))\\
\geq& \sum_{j \in A \cap B(i_0, \eps\theta)}  \beta(\alpha'_j, d(i_0, j)) + \sum_{j \in A \setminus B(i_0, \eps\theta)}  \beta(\alpha'_j, d(i_0, j)) + \sum_{j \in I}\beta((1-\delta)d(j,S), d(j, i_0))\,. 
\end{align*}
For the inequality above, for the first sum, the claim is trivial if $\alpha_j'$ is the minimizer, so let us assume that $\alpha_j' \geq (1+\eps)\theta - d(i_0, j)$. Note that $\alpha_j' \leq (1+\eps^2)\theta$.
Then, using that $d(i_0, j) \leq \epsilon \theta$ we have 
\begin{align*}
(1+\epsilon) \theta - d(i_0, j)-(\gamma (1-\delta/2)) d(i_0, j)
&= (1+\epsilon) \theta - (1-\eps)d(i_0, j)-(\gamma (1-\delta/2) + \eps) d(i_0, j)\\
&\geq (1+\eps^2) \theta - (\gamma (1-\delta/2) + \eps)d(i_0,j)\\
&\geq (1+\eps^2)\theta - \gamma d(i_0, j)\\
&\geq \alpha_j' - \gamma d(i_0, j)\,,
\end{align*}
where we used that $\gamma \geq \gamma (1-\delta/2) + \eps$ for $\delta \gg \eps$.

For the second sum, using that $d(i_0, j) > \epsilon \theta$ we have 
\[\theta - \gamma (1 - \delta/2) d(i_0, j) \geq (1+\eps^2) \theta - \gamma d(i_0, j) \geq \alpha_j' - \gamma d(i_0, j)\,,\]
where we used that $\gamma \delta d(i_0, j) / 2 \geq \gamma \delta \eps \theta \geq \eps^2 \theta$ for $\delta \gg \eps$ and $\alpha_j' \leq (1+\eps^2) \theta$.
The third sum follows trivially.
Then
\[\sum_{j \in A} \beta(\tilde\tau_j, (1-\delta/2) d(i_0, j)) + \sum_{j \in I}\beta((1-\delta)d(j,S), (1-\delta/2)d(j, i_0)) \geq \hat{f}\,.\]

We prove now that the other assumptions of the current lemma are also satisfied for $i_0$ and $\{\tilde\tau\}_{j \in A}$, thus reaching a contradiction because $\max_{j \in A}\tilde\tau_j \leq \tau'_{\max} < \tau_{\max}$.
We have by definition $\tilde\tau_j = \theta$ for $j \in A \setminus B(i_0, \eps \theta)$ and $\theta \leq \tilde\tau_j \leq \min((1-\delta)d(j,S),(1+\eps^2)\theta)$ for $j \in A \cap B(i_0, \eps \theta)$.

It remains to check the triangle inequalities $\tilde\tau_j \leq \tilde\tau_k + d(j, k)$.
The triangle inequalities between $j, k$ with $\tilde\tau_k \geq \tilde\tau_j$ are trivially satisfied, so we assume $\tilde\tau_j > \tilde\tau_k \geq \theta$.
Then $\tilde\tau_j = \min(\alpha'_j, (1+\epsilon) \theta - d(i_0, j))$, where $d(i_0, j) \leq \eps \theta$.
First, consider the case $d(i_0, k) > \eps \theta$, which implies $\tilde\tau_k = \theta$.
Then
\[\tilde\tau_j \leq (1+\eps)\theta - d(i_0, j) \leq (1+\eps)\theta - (d(i_0, k) - d(j, k)) \leq \theta + d(j, k) = \tilde\tau_k + d(j, k)\,.\]
Second, consider the case $d(i_0, k) \leq \eps \theta$, which implies $\tilde\tau_k = \min(\alpha'_k, (1+\epsilon) \theta - d(i_0, k))$, where also recall that $\alpha'_k = \min(\tau'_{\max}, (1-\delta)d(k, S))$.
If $\tilde\tau_k = \tau'_{\max}$, we are done.
If $\tilde\tau_k = (1-\delta)d(k, S)$, we have 
\[\tilde\tau_j \leq (1-\delta) d(j,S) \leq (1-\delta)(d(j,k)+d(k,S)) \leq \tilde\tau_k + d(j,k)\,.\]
Finally, if $\tilde\tau_k = (1+\epsilon) \theta - d(i_0, k)$, we have
\[\tilde\tau_j \leq (1+\epsilon) \theta - d(i_0, j) \leq (1+\epsilon) \theta - (d(i_0, k) - d(j, k)) = \tilde\tau_k + d(j,k)\,.\]
Hence, we satisfy all assumptions of the current lemma and we reach a contradiction.

\textbf{Fourth constraint for $i_0 \in S$.}
Finally, the remaining possibility is that there exists some $i_0 \in S$ for which the fourth constraint is violated.
However, in that case, we note that $\tau_j \leq (1-\delta) d(j, S) \leq d(i_0, j)$, so the left-hand side of the fourth constraint is zero, so it is impossible for the constraint to be violated.
\end{proof}

Given Lemmas~\ref{lem:pay-to-open-1} and~\ref{lem:pay-to-open-2}, the analysis of the merging procedure is nearly identical to that in~\cite{MR4928455-Cohen25}.
To prove that the merging procedure maintains the maximality of the sequences produced they (implicitly) use the fact that a ``payable'' facility is an ``openable'' facility.
This property holds in our case by Lemmas~\ref{lem:pay-to-open-1} and~\ref{lem:pay-to-open-2}.
The only remaining claim whose proof we need to change is that of Claim $4$ in~\cite{MR4928455-Cohen25}, which proves that if a sequence $H'_t$ is $(\nu, 0)$-valid with respect to some $\hat f', u'$, then it is also $(\nu, \eta)$-valid with respect to $\hat f, u$ that are $\eta$-close to $\hat f', u$.
We prove below a version of this claim that applies to our case, by starting with a verbatim copy of their proof and adapting it to our setting.
(This proof will also use Lemmas~\ref{lem:pay-to-open-1} and~\ref{lem:pay-to-open-2}.)
The rest of the analysis is identical up to cosmetic changes.

\begin{claim}[See Claim 4 in~\cite{MR4928455-Cohen25}]
\label{claim:param-ident}
    Let $H'_t$ be a $(\nu, 0)$-valid sequence with respect to some $\hat f', u'$.
    Let $\hat f, u$ satisfy that either $\hat f, u$ and $\hat f', u'$ are identical, or $\hat f > \hat f'$ with $|\hat f - \hat f'| \leq \eta$, or there exists a single free facility $\tilde{i}$ with $u(\tilde i) < u'(\tilde i)$  and $|u(\tilde i) - u'(\tilde i)| \leq \eta$.
    Then $H'_t$ is a $(\nu,\eta)$-valid sequence with respect to $\hat f, u$.
    \label{claim:1ststep_etavalid}
\end{claim}
\begin{proof}
Suppose first that $\hat f\neq \hat f'$ is the difference parameter.
Because each facility in $H'_p$ was $(\nu(p),0)$-openable with respect to parameters $\hat f', u'$, then it is also $(\nu(p),\eta)$-openable with respect to parameters $\hat f, u$ because the only difference is that $\hat f$ is larger than $\hat f'$ by at most $\eta$.
Furthermore, $H'_p$ is maximal with respect to $\hat f, u$ by the following argument based on Lemmas~\ref{lem:pay-to-open-1} and~\ref{lem:pay-to-open-2}: Since $H'_p$ was maximal with respect to $\hat f', u'$, it means the assumptions of Lemmas~\ref{lem:pay-to-open-1} (if $\nu(p)=1$) or~\ref{lem:pay-to-open-2} (if $\nu(p)=2$) are not satisfied with respect to $\hat f', u'$.
We know that their first assumption is always satisfied by Lemma~\ref{lem:seq-overbidding}, so the second assumption must be the one not satisfied.
However, their second assumptions become even stronger for $\hat f$ larger than $\hat f'$, so they are still not satisfied, which implies that there still does not exist any $(\nu(p),0)$-openable facility, because otherwise these assumptions would trivially be satisfied.

Next, consider the case when $u(\tilde i) \neq u'(\tilde i)$ is the difference parameter.
Then, by the same argument as above, we also have that each facility in $H'_p$ is $(\nu(p),\eta)$-openable with respect to parameters $\hat f, d$.
To see this more formally, let $H'_p = \langle h_1, \dots, h_{\ell} \rangle$ and consider $h_t$ for some $t \in [\ell]$, which is $(\nu(p), 0)$-openable with respect to $\hat f', u', \theta = (1 + \eps^2)^{\lfloor \frac{p - 1}{2} \rfloor}$ and some superset $S \supseteq (\cup_{q \in [p-1]} H_q) \cup \{ h_1, \dots, h_{t-1} \}$.
When we switch from $\hat f', u'$ to $\hat f, u$, because the only difference is $u(\tilde i) < u'(\tilde i)$, the set of active clients $A'$ for $\hat f', u'$ is a superset of the set of active clients $A$ for $\hat f, u$.
Also, let $d'(\cdot, \cdot)$ and $d(\cdot, \cdot)$ be the distances with respect to $u'$ and $u$, respectively. 
For any $h_t \in H_p$, we have that $h_t$ was $(\nu(p),0)$-openable with respect to the former setting with bids $(\tau'_j)_{j \in A'}$.
Then let $(\tau_j)_{j \in A'}$ be the new bids defined as $\tau_j = \min(\tau'_j, (1-\delta)d(j, S))$ for $j \in A'$ (these bids are only used on $A$, but we define them on the superset $A'$ for the purposes of the analysis). 
Let us verify that $h_t$ is $(\nu(p),\eta)$-openable with respect to $\hat f, u$ with the bids $(\tau_j)_{j \in A}$.
We need to check the corresponding conditions in Definitions~\ref{def:openable-mod-eta-1} (if $\nu(p)=1$) and~\ref{def:openable-mod-eta-2} (if $\nu(p)=2$).
We focus for simplicity on the case $\nu(p)=2$, which is the more intricate of the two.
The case $\nu(p)=1$ follows by similar arguments.

\begin{itemize}
    \item $\tau_j = \theta$ for every $j \in A - B(h_t, \eps \theta)$: Since it was true for $A' \supseteq A$, it remains true. 
    \item $\theta \leq \tau_j \leq \min\{(1-\delta)d(j,S),(1+\eps^2)\theta\}$ for every $j \in A \cap B(h_t, \eps \theta)$: Fix $j \in A \cap B(h_t, \eps \theta)$. 
    Since $A \subseteq A'$, $\tau'_j \leq \min\{(1-\delta)d'(j,S),(1+\eps^2)\theta\}$. Since $\tau_j$ is defined as $\min(\tau'_j, (1-\delta)d(j,S))$, the condition is satisfied.
    \item $
    \sum_{j \in A} \beta(\tau_j, (1-\delta/2)d(h_t,j)) + \sum_{j \in I} \beta((1-\delta)d(j,S), (1-\delta/2)d(h_t,j)) \geq \hat f - n \eta$: Let us argue that the contribution of each client $j$ to the left-hand side changes by at most $\eta$ between $(A', d', \tau')$ and $(A, d, \tau)$.
    We know that 
    \[\sum_{j \in A'} \beta(\tau'_j, (1-\delta/2)d(h_t,j)) + \sum_{j \in I'} \beta((1-\delta)d'(j,S), (1-\delta/2)d(h_t,j)) \geq \hat f\,.\]
    For $j \in A$, note that $\tau_j < \tau'_j$ implies $\tau_j = (1-\delta)d(j, S) \geq (1-\delta)(d'(j, S) - \eta) \geq \tau'_j - \eta$, so $\tau_j \geq \tau'_j - \eta$ in any case.
    Then 
    \[\sum_{j \in A'} \beta(\tau_j + \eta, (1-\delta/2)d(h_t,j)) + \sum_{j \in I'} \beta((1-\delta)d'(j,S), (1-\delta/2)d(h_t,j)) \geq \hat f\,.\]
    For $j \in I'$, note that $d(j, S)$ and $d'(j, S)$ only differ by at most $\eta$.
    So also 
    \[\sum_{j \in A'} \beta(\tau_j + \eta, (1-\delta/2)d(h_t,j)) + \sum_{j \in I'} \beta((1-\delta)(d(j,S) + \eta), (1-\delta/2)d(h_t,j)) \geq \hat f\,.\]
    Finally, for $j \in A' \cap I$, $(1-\delta)d(j, S) \geq \tau_j$, so $\beta(\tau_j + \eta, (1-\delta/2)d(h_t, j)) \leq \beta((1-\delta)d(j,S)+\eta, (1-\delta/2)d(h_t, j))$, so 
    \[\sum_{j \in A} \beta(\tau_j + \eta, (1-\delta/2)d(h_t,j)) + \sum_{j \in I} \beta((1-\delta) d(j,S) + \eta, (1-\delta/2) d(h_t,j)) \geq \hat f\,,\]
    from which it follows that 
    \[\sum_{j \in A} \beta(\tau_j, (1-\delta/2)d(h_t,j)) +  \sum_{j \in I} \beta((1-\delta)d(j,S), (1-\delta/2)d(h_t,j)) \geq \hat f - n \eta\,.\]
    \item For every regular facility $i_0$ and $k \in A$, $
    \sum_{j \in A} \beta(\tau_j, d(i_0,j)) + \sum_{j \in I} \beta((1-\delta)d(j,S),d(i_0,j)) \leq \hat f$: We have
    \[
    \sum_{j \in A'} \beta(\tau'_j, d'(i_0,j)) + \sum_{j \in I'} \beta((1-\delta)d'(j,S), d'(i_0,j))
    \leq \hat f.
    \]
    We have trivially that $d'(j, S) \geq d(j, S)$ and $\tau_j' \geq \tau_j$.
    Because $i_0$ is a regular facility, we also have $d'(i_0, j) = d(i_0, j)$.
    Also, for $j \in A' \cap I$, because $j \in I$ implies that $\tau_j' \geq \theta \geq (1-\delta)d(j, S)$, we have $\beta(\tau'_j, d(i_0, j)) \geq \beta((1-\delta) d(j,S), d(i_0, j))$, so we also have
    \[
    \sum_{j \in A} \beta(\tau_j, d(i_0,j)) + \sum_{j \in I} \beta((1-\delta)d(j,S), d(i_0,j))
    \leq \hat f.
    \]
    \item $\tau_j \leq \tau_k + d(j, k)$ for every $j, k \in A$: For $j, k \in A'$, we had $\tau'_j \leq \tau'_k + d(j,k)$.
    Recall that $\tau_j = \min(\tau_j', (1-\delta)d(j,S))$ for all $j \in A$.
    Suppose $\tau_k < \tau'_k$, as otherwise the inequality holds trivially.
    Then $\tau_k = (1-\delta) d(k, S)$, and also $\tau_j \leq (1-\delta) d(j, S)$, so
    \[\tau_j \leq (1-\delta) (d(j, k) + d(k, S)) = \tau_k + (1-\delta) d(j, k) \leq \tau_k + d(j, k)\,.\]
\end{itemize}
Furthermore, $H'_p$ is maximal w.r.t. $\hat f, u$ by the following argument based on Lemmas~\ref{lem:pay-to-open-1} and~\ref{lem:pay-to-open-2}: Since $H'_p$ was maximal w.r.t. $\hat f', u'$, it means the assumptions of Lemmas~\ref{lem:pay-to-open-1} (if $\nu(p)=1$) or~\ref{lem:pay-to-open-2} (if $\nu(p)=2$) are not satisfied w.r.t. $\hat f', u'$.
We know that their first assumption (about $\leq \hat{f}$) is always satisfied by Lemma~\ref{lem:seq-overbidding}, so the second assumption must be the one not satisfied.
Let us argue that these second assumptions are still not satisfied w.r.t. $\hat f, u$:
\begin{itemize}
    \item If $\nu(p)=1$, for Lemma~\ref{lem:pay-to-open-1}, we know that for all regular facilities $i \in F \setminus S$
    \[\sum_{j \in A'} \beta(\theta, (1-\delta)d(i,j)) + \sum_{j \in I'} \beta((1-\delta)d'(j, S), (1-\delta)d(j,i)) < \hat{f} - n \eta\,.\]
    We have $d'(j, S) \geq d(j, S)$, and also for $j \in A' \cap I$ we have that $\theta \geq (1-\delta) d(j, S)$, so it easily follows that 
    \[\sum_{j \in A} \beta(\theta, (1-\delta)d(i,j)) + \sum_{j \in I} \beta((1-\delta)d(j, S), (1-\delta)d(j,i)) < \hat{f} - n \eta\,.\]
    \item If $\nu(p)=2$, for Lemma~\ref{lem:pay-to-open-2}, we know that for all regular facilities $i \in F \setminus S$ and all $\{\tilde{\tau}_j\}_{j \in A'}$ such that $\tilde{\tau}_j = \theta$ for $j \in A' \setminus B(i, \epsilon \theta)$, $\theta \leq \tilde{\tau}_j \leq \min((1-\delta)d'(j,S), (1+\epsilon^2)\theta)$ for $j \in A' \cap B(i, \epsilon \theta)$, $\tilde{\tau}_j \leq \tilde{\tau}_k + d(j, k)$ for all $j, k \in A'$, we have
    \[\sum_{j \in A'} \beta(\tilde\tau_j, (1-\delta/2)d(i,j)) + \sum_{j \in I'} \beta((1-\delta)d'(j, S), (1-\delta/2)d(j,i)) < \hat{f} - n \eta\,.\]
    We will proceed by contradiction. 
    Suppose there exists a regular facility $i \in F \setminus S$ and some $\{\tilde{\tau}_j\}_{j \in A}$ such that $\tilde{\tau}_j = \theta$ for $j \in A \setminus B(i, \epsilon \theta)$, $\theta \leq \tilde{\tau}_j \leq \min((1-\delta)d(j,S), (1+\epsilon^2)\theta)$ for $j \in A \cap B(i, \epsilon \theta)$, and $\tilde{\tau}_j \leq \tilde{\tau}_k + d(j, k)$ for all $j, k \in A$, with
    \[\sum_{j \in A} \beta(\tilde\tau_j, (1-\delta/2)d(i,j)) + \sum_{j \in I} \beta((1-\delta)d(j, S), (1-\delta/2)d(j,i)) \geq \hat{f} - n \eta\,.\]
    Let us extend $\tilde\tau_j = \theta$ for all $j \in A' \setminus A$, which trivially satisfies $\tilde\tau_j \leq \min((1-\delta) d'(j, S), (1+\eps^2)\theta)$ for $j \in A' \cap B(i, \eps \theta)$.
    Note also that $d(j, S) \leq d'(j, S)$.
    Furthermore, for $j \in A' \setminus A$, we have $\theta \geq (1-\delta)d(j,S)$, so
    \[\sum_{j \in A'} \beta(\tilde\tau_j, (1-\delta/2)d(i,j)) + \sum_{j \in I'} \beta((1-\delta)d'(j, S), (1-\delta/2)d(j,i)) \geq \hat{f} - n \eta\,.\]
    To reach a contradiction, it remains to prove that $\tilde\tau_j \leq \tilde\tau_k + d(j,k)$ for all $j, k \in A'$.
    We already know this for $j, k \in A$, and the claim is also trivial if $j \in A' \setminus A$ because $\tilde\tau_j = \theta \leq \tau_k$.
    So let us consider the case $j \in A$ and $k \in A' \setminus A$.
    Note that in this case $\tilde\tau_j \leq (1-\delta) d(j,S)$ and $\tilde\tau_k = \theta \geq (1-\delta) d(k, S)$.
    Then $\tilde\tau_j \leq (1-\delta) (d(j,k)+d(k,S)) \leq \tilde\tau_k + d(j,k)$, so we reach a contradiction.
\end{itemize}
Therefore, Lemmas~\ref{lem:pay-to-open-1} and~\ref{lem:pay-to-open-2} are still not satisfied, which implies that there still does not exist any $(\nu(p),\eta)$-openable facility, because otherwise these assumptions would trivially be satisfied.
\end{proof}